\appendix
\renewcommand{\thesection}{\Alph{section}}
\begin{center}{\Large\bfseries 附录概览}\end{center}
本附录提供支撑正文的实现细节、实验协议、训练规格和扩展分析。

\vspace{4pt}
{\small
\newcommand{\tocA}[3]{\textbf{#1}\quad\textbf{#2}\dotfill\pageref{#3}\par\vspace{-1pt}}
\newcommand{\tocD}[1]{\hspace*{1.6em}{\color{black!65}#1}\par\vspace{2pt}}
\newcommand{\tocB}[3]{\hspace*{1.6em}#1\quad #2\dotfill\pageref{#3}\par\vspace{-1pt}}
\tocA{A}{可复现性索引}{app:repro}
\tocD{产物、发布状态，以及各项协议细节的说明位置。}
\tocA{B}{局限性}{app:limit}
\tocD{自主流程的适用范围、所评估的维度以及项目规模。}
\tocA{C}{\rsig{} 的系统设计与实现}{app:system}
\tocD{递归开发系统及其智能体接口的完整规格。}
\tocB{C.1}{递归开发流程}{app:proc}
\tocB{C.2}{全局质量监控}{app:monitor}
\tocB{C.3}{高层导演界面}{app:director}
\tocB{C.4}{智能体提示词与输出约定}{app:prompts}
\vspace{3pt}
\tocA{D}{实验协议与评估可靠性}{app:exp}
\tocD{所报告实验背后的配置、核算与可靠性协议。}
\tocB{D.1}{实验配置}{app:config}
\tocB{D.2}{Token 与成本核算}{app:cost}
\tocB{D.3}{回放与评分稳定性}{app:stability}
\tocB{D.4}{自由游玩评估}{app:freeplay}
\tocB{D.5}{主要结果的统计可靠性}{app:stat}
\tocB{D.6}{自适应方向策略评估}{app:adaptive}
\tocB{D.7}{智能体验证审计}{app:audit}
\tocB{D.8}{多阶段指引评估}{app:multistage}
\vspace{3pt}
\tocA{E}{训练数据与经验内化}{app:train}
\tocD{游戏开发经验的构建、整理与内化。}
\tocB{E.1}{训练语料的构建与整理}{app:corpus}
\tocB{E.2}{微调配置}{app:ft}
\tocB{E.3}{训练数据消融}{app:ablation}
\vspace{3pt}
\tocA{F}{扩展评估与案例研究}{app:extended}
\tocD{扩展比较、定性轨迹以及按任务类别的细分结果。}
\tocB{F.1}{与多智能体“生成–验证”系统的比较}{app:vibegame}
\tocB{F.2}{定性演化案例研究}{app:cases}
\tocB{F.3}{按类别的结果}{app:family}
}
\clearpage

\section{可复现性索引}\label{app:repro}
表~\ref{tab:repro} 即可复现性声明中提到的统一索引：对于每项产物，它给出该产物包含的内容、是否已发布，以及相应协议或实现细节的说明位置。

\begin{table}[htbp]
\caption{\textbf{可复现性索引。}产物、发布状态，以及相应实现与协议细节的位置。}
\label{tab:repro}
\centering\small
\begin{tabular}{@{}p{0.27\linewidth}p{0.47\linewidth}p{0.2\linewidth}@{}}
\toprule
条目 & 内容 / 规格 & 位置 \\
\midrule
\multicolumn{3}{@{}l}{\textit{\kai 已发布的产物}} \\
实现 & 局部/全局循环、全局质量监控器、导演评审工具和评分框架 & 代码，已上线 \\
提示词与输出约定 & 所有智能体使用的完整提示词和结构化输出 & 代码，已上线；附录~\ref{app:prompts} \\
可游玩的演化过程 & 六个具代表性的游戏及其保留的检查点 & 项目主页，已上线 \\
冻结的基础项目 $P_0$ & 每个任务、每个生成器的起始项目 & 待发布 \\
评估产物 & 逐任务的细则分数、评审输出和回放报告（51,644 个文件） & HF，已上线$^\dagger$ \\
训练语料 & 2,213 条生成轨迹、2,108 份规划和 2,013 个经验证的改进轮次 & 待发布 \\
运行树 & 长周期开发运行的逐轮项目快照 & 待发布 \\
\midrule
\multicolumn{3}{@{}l}{\textit{\kai 协议与实现细节}} \\
递归开发循环 & 一个局部轮次、全局检查点更新和停止逻辑 & 附录~\ref{app:proc} \\
全局质量监控器 & 最佳检查点选择、饱和准则和耐心值 $K$ & 附录~\ref{app:monitor} \\
模型分配 & 用于生成、探索、编辑、验证和评审的模型 & 表~\ref{tab:models} \\
超参数 & 循环、监控器和评分参数 & 表~\ref{tab:params} \\
开发预算 & 每个报告行的轮次与工具调用预算 & 附录~\ref{app:config} \\
Token 与成本核算 & 计费 token、排除的成本和定价 & 附录~\ref{app:cost} \\
评分稳定性 & 回放与评审的变异性 & 附录~\ref{app:stability} \\
语料构建 & 抽取、验证过滤和污染检查 & 附录~\ref{app:corpus} \\
微调 & 训练数据混合与超参数 & 附录~\ref{app:ft} \\
导演评审 & 导演的输入、交互预算和指引格式 & 附录~\ref{app:director} \\
按类别的结果 & 按任务类别细分的主表结果 & 附录~\ref{app:family} \\
\bottomrule
\end{tabular}

\vspace{3pt}
\parbox{0.92\linewidth}{\footnotesize $^\dagger$ 发布于 Hugging Face：\texttt{RSIGame/RSIGame-TableArtifacts}。发布前已移除机器特定路径和可识别身份的术语，导出的文件经过重新解析和重新扫描。}
\end{table}

\section{局限性}\label{app:limit}
\rsig{} 的核心流程是完全自主的：自适应的局部开发、全局质量监控、最佳检查点跟踪和感知饱和的停止都不需要人工干预。我们另外研究了一种可选的有引导扩展：只有在自主改进饱和之后，才引入来自人类或更强模型的稀疏高层指引。目前的研究只在有限的多阶段开发案例上评估了这一机制；关于何时应调用指引、多少指引是有益的、以及人类指引与模型指引有何差异，更广泛的研究留待未来工作。

我们的评估也集中于机制、深度、视觉和美术四个维度。这些维度并不能完全刻画更高层次的属性，例如原创性、叙事质量、长期玩家参与度或主观乐趣。

最后，我们的实验关注的是相对紧凑、能够在实际的智能体预算内迭代开发的游戏。将自主递归开发扩展到规模大得多、周期更长、素材更丰富且跨系统依赖更复杂的项目，仍然是一个重要的方向。

\section{\rsig{} 的系统设计与实现}\label{app:system}
\subsection{递归开发流程}\label{app:proc}
算法~\ref{alg:rsig} 是第~3 节循环的完整写法：内循环的一个轮次、全局质量监控器运行的检查点，以及饱和可能导致的两种结果——返回冠军版本，或接受一条指引并开启下一阶段。

多阶段结构由计数器 $c$ 和阶段索引 $s$ 体现。一个检查点要么晋升出新的冠军（将 $c$ 归零），要么没有找到更好的版本（$c$ 加一）；连续 $K$ 个检查点都没有晋升，就是本文所说的饱和。饱和本身并不结束开发：它结束的是当前阶段。如果有可用的指引，就针对冠军写一份简报，下一阶段从该冠军（而不是上一轮的构建）重新开始，$c$ 被重置，循环在 $\gamma_{s+1}$ 引导后续每一轮的情况下继续。只有当某个阶段饱和且不再给出任何指引时，开发才会结束。本文的运行使用 $E = 3$ 和 $K = 3$。

\begin{algorithm}[h]
\caption{\rsig{} 递归游戏开发}
\label{alg:rsig}
\small
\begin{algorithmic}[1]
\Require 游戏规格说明 $x$，初始项目 $P_0$，初始清单 $C_0$，检查点间隔 $E$，耐心值 $K$
\State $P^\star_0 \gets P_0,\ \gamma_0 \gets \varnothing,\ t \gets 0,\ s \gets 0,\ c \gets 0$ \Comment{冠军、指引、轮次、阶段、空闲检查点数}
\While{仍有开发预算}
  \State $d_t \sim \pi_{\mathrm{controller}}(\cdot \mid x, P_t, C_t, \gamma_s)$ \Comment{清单和指引引导本轮}
  \State $\tau_t \sim \pi_{\mathrm{explorer}}(\cdot \mid P_t, d_t)$
  \State $\Delta_t \sim \pi_{\mathrm{editor}}(\cdot \mid P_t, d_t, \tau_t),\ P_{t+1} \gets \mathrm{Apply}(P_t, \Delta_t)$
  \State $v_t \sim \pi_{\mathrm{verifier}}(\cdot \mid P_{t+1}, d_t, \tau_t)$ \Comment{改进是否可观察，什么发生了退化}
  \State $C_{t+1} \gets U(C_t, d_t, \tau_t, \Delta_t, v_t)$ \Comment{下一轮赖以构建的工作记忆}
  \If{$t \bmod E = 0$} \Comment{检查点：由监控器比较，而非循环本身}
    \State $P' \gets \mathrm{SelectBest}(P^\star_s, P_{t+1})$
    \If{$P' \neq P^\star_s$}
      \State $P^\star_s \gets P';\ c \gets 0$ \Comment{新冠军；本阶段仍有回报}
    \Else
      \State $c \gets c + 1$ \Comment{又一个没有更好版本的检查点}
    \EndIf
    \If{$c \geq K$} \Comment{饱和：$K$ 个检查点无新冠军}
      \If{没有可用的指引} \Return $P^\star_s$ \EndIf
      \State $\gamma_{s+1} \gets \textsc{Guidance}(x, P^\star_s)$ \Comment{针对冠军的简报开启下一阶段}
      \State $P^\star_{s+1} \gets P^\star_s,\ P_{t+1} \gets P^\star_s,\ s \gets s+1,\ c \gets 0$ \Comment{下一阶段从这里重新开始}
    \EndIf
  \EndIf
  \State $t \gets t + 1$
\EndWhile
\State \Return $P^\star_s$
\end{algorithmic}
\end{algorithm}

\subsection{全局质量监控}\label{app:monitor}
\paragraph{比较什么。}监控器从不单独给某个版本打分。在每个检查点，它在同一组演示上回放保留的冠军 $P^\star_s$ 和候选构建，逐个场景地将两段录像配对，并且只问两者之间发生了什么变化。比较器不会被告知候选来自哪一轮、该改进试图实现什么，也不会被告知已花费了多少预算：方向和预算在上一层决定，如果在这里提供它们，开发意图就可能影响模型所报告看到的内容。

\paragraph{量表。}代理评分细则的每条准则都在三个锚定等级上读取——1.0 表示按描述工作，0.5 表示部分工作、间歇工作或仅在某些情况下工作，0.0 表示不工作——冠军和候选分别读取，并且只针对两段录像都展示到的准则，从而使差异反映的是游戏本身的差异，而不是两次回放恰好到达的内容不同。各准则读数被合并为 $[0,1]$ 区间内的加权总值，并分为两半报告：来自当前阶段简报的准则，以及来自任务文档的准则。

\paragraph{采纳。}当简报那一半提高至少 0.05、且任务那一半下降不超过 0.05 时，候选取代冠军；否则保留冠军。任务那一半起到的是底线而非目标的作用，因此只要构建没有让任务已要求的内容退化，就可以因为推进了本阶段自身的目标而被采纳。如果没有阶段准则，该规则就退化为总值提高 0.05。

\paragraph{检查点与停止。}检查点每三轮进行一次，这是能够包含一次改进、其验证和一次回放的最短间隔，因此也是能观察到差异的最短窗口。一旦冠军连续 $K$ 个检查点保持不变，开发即停止；我们使用 $K = 3$。图~\ref{fig:global}(b) 在相同的运行上将 $K = 1 \ldots 4$ 与固定预算进行比较：较小的 $K$ 停得更早、交付得更少，较大的 $K$ 以固定预算轮次的一小部分接近其得分，而 $K = 3$ 恰好位于拐点。

\paragraph{与评估器隔离。}监控器与基准评审只共享回放机制。监控器读取的是一份单独的代理评分细则，从不读取基准的保留细则，也从不读取其分数或反馈；该模块完全不导入评审。这种隔离在内部同样成立：比较器是唯一看到像素的组件，而决定方向和停止的阶段读取的是它的结构化计数而非画面帧，因此没有任何一层既能解读图像，又能依据自己的解读采取行动。

\subsection{高层导演界面}\label{app:director}
第~3 节描述了阶段简报对循环的作用。本附录描述简报从何而来。本文报告的每一份简报——无论来自人类还是模型——都是在下述界面中、按同一协议写成的，因此“由人指导这次运行”和“由模型指导这次运行”之间，除了导演本身之外别无差异。

\subsubsection{导演获得什么}
导演评审的是某个游戏的某一个构建版本。图~\ref{fig:director} 就是该页面：游戏在左侧面板中实时运行，框架获得焦点时接管键盘；右栏显示预算、生成该游戏所依据的任务规格说明，以及已发生开发过程的摘要；下方的表单就是简报。

这份摘要刻意保持简略。它报告已运行了多少个自主轮次、其中哪些在全局质量监控器的每三轮检查点上成为冠军、正在评审的是哪个构建，以及这些轮次如何在功能性工作和呈现性工作之间分配。它不报告任何构建在任何一轮的任何分数。

这种省略正是关键所在。导演被问到的是游戏接下来需要什么，而分数回答的是另一个问题——上一轮是否奏效——而且用的是循环本身已经在优化的那种“货币”。关于运行的其他任何信息也都不会出现在页面上：没有源代码、没有日志、没有改进过程记录、没有其他导演的评审。这个工具展示的只是一个可玩的游戏，以及它本应成为的那份文档。

\begin{figure}[t]
\centering
\fig{8}
\caption{导演评审页面（会话进行中），展示的是已经历六个监控器检查点的 \textit{Alley Brawlers} 构建：实时游戏、预算、任务规格说明和开发摘要——冠军出现在 R03 和 R09，此后再无更新，以及十八轮在功能性与呈现性工作之间的分配。画面是一个真实的游戏瞬间，展示了两位导演随后用不同措辞描述的缺陷：玩家的格斗者已经散成一条杂乱的精灵图块带。其下的简报表单见图~\ref{fig:brief}。}
\label{fig:director}
\end{figure}

\subsubsection{同一工具，同一预算}
导演有 60 次操作和 2 次重置。按一次键无论按住多久都算一次操作，点击一次也算一次，而观看和等待是免费的，因此预算买的是交互而不是耐心。图~\ref{fig:director} 中的计数器就是会话实际执行的计数器，每位导演实际花费的次数随简报一起保存。

人类通过浏览器游玩。模型则通过一个命令行客户端游玩——\texttt{open}、\texttt{press}、\texttt{hold}、\texttt{click}、\texttt{wait}、\texttt{look}、\texttt{reset}、\texttt{state}、\texttt{submit}——它驱动同一服务器上的同一个会话对象，通过同样的计数器花费同样的预算，并在每次操作后返回捕获的画面帧供模型阅读。两位导演都看不到对方的传输方式；二者面对的是同一个正在运行的游戏。

两位导演评审的是完全相同的字节。每份简报都记录了其所针对的项目树的哈希值，对于 \textit{Alley Brawlers}，两者读到的都是 \texttt{c12fe90b30b6c176}。人类导演用满了 60 次操作，没有重置；模型用了 25 次，没有重置。

\subsubsection{模型如何担任导演}
模型导演获得一份冻结的、带版本号的提示词（\texttt{director-prompt/2}），在一次运行中保持不变；修改它就构成一次新的运行。该提示词让模型扮演一位资深游戏导演，评审一个自主开发已陷入停滞的构建；要求它决定下一阶段应该是什么，而不是罗列缺陷；并禁止给出实现层面的指令。

有两条规则承担了方法论上的作用。模型对游戏的一切了解都必须来自会话本身：读取、列出或搜索任何其他文件——任何源代码、日志、分数或其他评审——都会使该评审作废。并且简报必须通过表单提交所用的同一个端点、以同一种格式提交，且只能提交一次；提交即结束会话。

\subsubsection{简报是什么}
四个字段，对两位导演完全相同：一句话的 \texttt{stage\_objective}（阶段目标）；\texttt{why\_now}（为何是现在），以导演在游玩中所见为依据；一到三条 \texttt{priorities}（优先事项），以方向而非代码修改的形式表述；以及零到三条需要 \texttt{preserve}（保留）的内容。循环在三处使用简报——规划器、阶段排序，以及改进和美术提示词——并且监控器据以打分的阶段清单会被重建，使简报自身的准则占据最大权重（附录~\ref{app:monitor}）。

\begin{figure}[t]
\centering
\fig{9}
\caption{简报表单，即图~\ref{fig:director} 页面的下半部分。模型导演通过该表单提交所用的端点，以 JSON 形式提交同样的四个字段。}
\label{fig:brief}
\end{figure}

表~\ref{tab:briefs} 将 \textit{Alley Brawlers} 的两份简报并列展示。它们针对同一棵树、相隔四分钟写成，而且毫无重叠。人类把这个构建看作呈现问题，要求改进美术和新手引导；模型打了两场完整比赛，从未让对手掉过一点血，于是把这个构建看作机制问题。二者都不是唯一正确的答案，结果也很接近：从这个检查点出发，人类引导的分支达到 62.0，模型引导的分支达到 61.5，而从同一构建重启、不给任何简报的对照分支达到 60.0。这项比较表明的是，简报提供了一个循环自身的证据在十八轮中都未曾产生的方向，并且两位导演通过同一工具看同一个构建，可以给出截然不同的方向。

\begin{table}[t]
\centering\small
\begin{tabular}{@{}p{0.11\linewidth}p{0.85\linewidth}@{}}
\toprule
\multicolumn{2}{@{}l}{\textbf{人类导演}——60/60 次操作，0 次重置。} \\
\midrule
目标 & 美术需要改进：标题画面和游戏结束画面杂乱且没有像样的背景，角色美术很丑，应该用漫画风格重新生成，攻击也没有特效。游戏还需要教会玩家怎么玩——现在我根本不知道该按什么。 \\
为何是现在 & 角色能动了，也有了标题和游戏结束画面，但它仍然太粗糙，缺乏吸引力。 \\
\midrule
\multicolumn{2}{@{}l}{\textbf{模型导演（Claude Opus 5）}——25/60 次操作，0 次重置。} \\
\midrule
目标 & 把比赛从一场按剧本进行的单方面痛殴变成真正的双方对战：玩家的输入必须可见且可靠地命中、格挡和脱身，对手是一个可以被击败的 AI，两名格斗者在外观和打法上确实不同。 \\
为何是现在 & 在两场完整的比赛中（两个格斗者各用一次），AI 一个像素的血都没掉。每个回合都在大约 18 秒内结束，我被压在左边角落里陷入“命中、倒地、命中”的循环。轻攻击、重攻击、特殊指令输入、按住后退、下蹲和跳跃，全都无法换血、格挡或让我脱身。当按什么都不起作用时，没有人能学会玩这个游戏。[……] 七轮视觉润色都花在了打斗周围的框架上（标题、背景、横幅、结算画面），但打斗本身仍然缺席。 \\
\bottomrule
\end{tabular}
\captionsetup{position=bottom}
\caption{为同一个 \textit{Alley Brawlers} 构建（树 \texttt{c12fe90b30b6c176}）写下的两份简报，节选其前两个字段。限于篇幅，省略了优先事项和保留项。}
\label{tab:briefs}
\end{table}

\begin{figure}[t]
\centering
\fig{10}
\caption{导演打开的案例列表：六个游戏、各自正在评审的构建，以及剩余预算。案例标识符为内部标签。}
\label{fig:caselist}
\end{figure}

\subsubsection{案例集}
共有六个游戏接受了评审，每个游戏都在其自身运行收敛到的构建上接受评审，范围从刚生成的项目到第十八轮不等（图~\ref{fig:caselist}）。每个游戏都由人类导演和模型导演各自独立评审，随后每份简报都交给一个从同一构建重启的分支，并与一个不给简报的对照分支进行比较。六个游戏中有三个是附录~\ref{app:cases} 的案例研究。

\subsection{智能体提示词与输出约定}\label{app:prompts}
本附录报告承担循环决策的五个模块的提示词，每个模块对应第~3 节的一个阶段。每份提示词都被删节为其决策骨架：固定的目录、示例、画面标注约定和框架内务性内容都被删去，每处删减都标记为 \texttt{[...]}。诸如 \texttt{\{items\}} 的花括号是运行时插槽，由开发状态、试玩记录或回放画面填充。

有一个设计选择在全部五个模块中反复出现，统一说明一次更容易看清。每份提示词都被告知回答其自身问题所需的信息，而刻意不被告知其余信息。规划器得到的本轮方向是一个结论，而不是一个论证。验证器被告知本轮试图做什么，因为那正是它必须检验的主张；而比较器既不被告知有任何东西被改进过，也不被告知哪个构建更新。编辑器被告知问题已经确立，这样它就从“这在代码的哪里”开始，而不是重新推导诊断。正是这种“有所保留”，使每个阶段不会去确认自己的预期。

\begin{table}[t]
\caption{\rsig{} 循环中五个面向 LLM 的模块，按其实现的第~3 节阶段排列。}
\label{tab:modules}
\centering\footnotesize
\begin{tabular}{@{}p{0.15\linewidth}p{0.11\linewidth}p{0.36\linewidth}p{0.3\linewidth}@{}}
\toprule
阶段 & 模块 & 输出约定 & 使用方 \\
\midrule
方向决策 & 控制器 & 一个 \texttt{development\_question}；动作取自 \{\texttt{reuse}, \texttt{explore}, \texttt{escalate}\}；一份 \texttt{exploration\_plan}，指明能回答该问题的最低成本方式。 & 设定本轮目标；\texttt{escalate} 报告该方向已耗尽，而不是凭空编造工作。 \\
智能体探索 & 探索器 & 一次会话应解决的 1–3 条需求、要启动的场景，以及给游玩智能体的一条可执行指令。 & 驱动测试智能体与 $P_t$ 交互，产生轨迹 $\tau_t$。 \\
基于证据的编辑 & 编辑器 & 在构建约束下对项目的编辑，以及一份修改内容报告。 & 产生 $P_{t+1}$。 \\
智能体验证 & 验证器 & \texttt{goal\_achieved} 取自 \{\texttt{yes}, \texttt{no}, \texttt{unclear}\}；\texttt{broken\_now}；当该列表非空时给出 \texttt{next\_goal}。 & 接受或拒绝本轮；任何被破坏的内容都作为待办工作重新进入状态。 \\
全局质量监控 & 全局质量监控器 & 两个构建的逐准则读数、带方向和幅度的变化，以及退化情况——从不给出分数。 & 汇总为决定冠军的代理值（附录~\ref{app:monitor}）。 \\
\bottomrule
\end{tabular}
\end{table}

\subsubsection{控制器}
规划器决定的是要\emph{查明}什么，而从不决定要\emph{改变}什么。它得到的输入包括：作为既定结论的本轮方向、需求状态、构建上次被游玩时的表现、可用的演示以及剩余预算；它返回一个问题以及回答该问题的最低成本方式。提示词的大部分篇幅用于排除那些任何观察都无法解决的问题。

\begin{promptbox}{pbg1}{提示词 P1. 控制器}
\begin{PV}
你负责规划一个正在逐轮开发的游戏接下来要调查什么。

你不写代码、不提出补丁，也不说应该改什么。那是另一个
智能体的工作，在你要求的证据收集完之后进行。你的输出
是关于要“查明”什么的决定。

[... 方向、需求状态、上次试玩、演示、预算 ...]

====================
要决定的内容
====================
一个主要问题。它必须：
  - 位于上述方向之内；
  - 以上述状态和证据为依据，而不是以这类游戏通常
    具备什么为依据；
  - 能通过观看游戏来回答：有人回放它并观察。而不是
    通过阅读源代码——没有人会替你去读；
  - 足够具体，使你能说出那个观察结果会是什么；
  - 有助于判断某处是否需要修复。

不是问题：“游戏能被改进吗？”——没有任何观察能回答它。
不是问题：“改进 HUD 布局。”——那是一条修改指令，而且
目前还什么都没有观察到。

然后选择一个动作：
  reuse      上面的证据已经说明了问题是什么。
  explore    证据缺失或含糊，而某个具体观察能解决它。
  escalate   这个方向上已没有有依据的未决问题。
             如实说明，而不是编造一个。

如果选择探索，选出能回答该问题的最低成本方式：
  replay_existing_demo    已有演示恰好覆盖了你需要看的内容
  extend_existing_demo    某个演示能到达该状态，但在得出
                          答案前就停止了
  custom_targeted_probe   没有演示能到达。代价最高的选项。

如果某个问题已经连续问了好几轮仍未解决，用同样的方式
再问一次就不算计划。要么说明这次会有什么不同，要么
换一个问题。

只用严格的 JSON 回复，不要输出其他内容：
{"action": "reuse|explore|escalate",
 "development_question": {"question": "...", "why_now": "..."},
 "exploration_plan": {"mode": "...", "demo_id": "...", "focus": "...",
                      "stop_when": "..."}}
\end{PV}
\end{promptbox}

\subsubsection{探索器}
探索是有目标的，而非泛泛而为。探针选择器把本阶段尚无法回答的需求转化为一条游玩智能体可以执行的指令，并被要求在重要性与一次短会话能到达的范围之间进行权衡。

\begin{promptbox}{pbg2}{提示词 P2. 探索器}
\begin{PV}
你要决定下一次“单次”游玩会话应该去查明这个游戏的
什么。

以下是本阶段仍无法回答的需求，以及每条可以在哪里检查：

[... 候选需求，附状态和场景 ...]

选出“此刻”最重要的 1 到 {cap} 条，并为将要游玩该构建
的人写一条指令。判断重要性——游戏其余部分所依赖的需求
胜过细节，几步操作就能解决的胜过需要长时间刷的。

如果你选的几条可以在同一个地方观察，请说明；一次需要
启动两个不同状态的会话，会把一半步数花在路上。

指令必须“可执行”：说出要做的输入或要到达的画面，以及
要观察什么。“检查飞机是否会转向”太弱。“在一次出击中，
先按左再按右各约一秒，报告飞机的航向是否改变”才是可以
执行的。

只用 JSON 回复，不要输出其他内容：
{"items": ["T04"], "scenario": "", "ask": "...", "why": "一句话"}
\end{PV}
\end{promptbox}

\subsubsection{编辑器}
编辑器接收的是一个已经确立的问题，提示词的首要任务是阻止它再去确立一遍。它获得源代码、必须通过的构建命令，以及运行并观察构建的方法；它没有修改框架或评估轨迹的权限。

\begin{promptbox}{pbg3}{提示词 P3. 编辑器}
\begin{PV}
你在改进一个小游戏。这是第 {r} 轮，最多 {R} 轮。

你没有玩过这个构建。是别人玩的：一个游玩智能体驱动了
它，画面已被解读，下面是这些证据所支持的内容以及为本轮
选定的内容。

所以下面的问题已经确立。你不需要查明它是否真实，也不
应该重新回放游戏去确认——那项工作已经完成，重复它会
耗掉你修复所需的调用次数。你的工作从“这在代码的哪里”
开始。

[... 本轮的证据包、近期历史、项目地图 ...]

你拥有的
  - 源代码，就在你工作的目录中
  - `{build_cmd}` —— 你停止时它必须通过
  - 你可以运行构建并观察它：[...]

规则
  - 改游戏，不要改框架。不要编辑 `node_modules`、`dist`、
    `demo_outputs/`，或 `_repair_evidence/` 下的任何内容。
  - 不要编辑 `demo_outputs/`。那是评估器回放的输入轨迹；
    改动它们改变的是考卷，而不是游戏。
  [... 流程与服务器规范 ...]
\end{PV}
\end{promptbox}

\subsubsection{验证器}
编辑之后，每个随附的演示都会被回放，由一个阅读者判定本轮是否配得上提交。提示词区分了两种容易混淆的失败：游戏从未有过的功能不是缺陷；一个产生了可见效果的输入已经被接收，即使脚本始终没有到达它原本设定的结果。只有毫无改变的输入才算被破坏。

\begin{promptbox}{pbg4}{提示词 P4. 验证器}
\begin{PV}
一个修复智能体刚刚编辑了一个游戏。你的工作是看看游戏
现在的表现并回答两个问题。你是这次编辑唯一的检查：
在你之后，它就发布了。

本轮试图做什么
{goal}

[... 该目标所源自的需求；回放画面 ...]

按以下顺序回答这两个问题。

1. goal_achieved —— yes / no / unclear。
   本轮打算做的事情现在是否明显成立？“unclear”是一个
   真实的答案：当演示从未到达该目标涉及的那部分游戏时，
   就这样回答，而不要猜测。

2. broken_now —— 一个列表，可以为空。
   观看这些画面的玩家会明显觉得不对的任何表现。具体说出
   你“看到”了什么、在哪一帧。最重要的几类，因为它们在
   源代码里看不出来：
     - 输入完全没有引起变化——前后画面一次又一次相同，
       说明游戏没有在接收输入
     - 有东西盖住了游戏区域：某个精灵、特效或面板挡住了
       角色、HUD 或动作
     - 画面冻结：最后一帧与第一帧相同
     - 文字无法阅读、被截断或相互重叠
   不要在这里列出缺失的功能。游戏从未有过的功能属于任务
   需求，而不是被破坏的东西。

还有一件不要报告的事。一个明显“产生了”效果的输入——
出现了模式标签、显示了提示——已经被接收，即使演示始终
没有到达它原本设定的结果。那是游戏需要的步骤比这个固定
脚本执行的更多，并不是游戏的错误。只有“毫无”改变的
输入才属于 broken_now。

然后，仅当 broken_now 非空时，写下 next_goal：用一句话
告诉修复智能体首先修什么。修最严重的那个，而不是全部。
\end{PV}
\end{promptbox}

\subsubsection{全局质量监控器}
比较器是整个方法中唯一读取像素的组件。它看到的是在同一固定脚本下回放的冠军和候选，并被问到发生了什么变化；它既不被告知哪个构建更新，也不被告知发生过改进。提示词的大部分内容用于防范成对比较出错的两种方式：假定更新的构建更好，以及把录制时序误认为行为变化。

\begin{promptbox}{pbg5}{提示词 P5. 全局质量监控器}
\begin{PV}
同一个游戏的两个构建，用“相同”的固定输入脚本回放。
你的工作是说出这个场景中发生了什么变化，别无其他。

[... 场景、它能涉及的准则、配对的画面 ...]

脚本是固定的，所以两个构建在相同时刻收到相同的输入。
对于下面的每个时刻，你会先看到旧构建在输入前后的画面，
再看到新构建在同样两个时刻的画面。

录制并非逐帧精确：同一构建回放两次，截取同一时刻也可能
相差一两帧，仅这一点就会让分数读数差一分、动画看起来
慢半拍，或让下落的方块低一行。那是录制器，而不是构建。

正是出于这个原因，每次输入后你会得到两帧。到第二帧就
消失的差异，或被之后的时刻抵消的差异，属于时序。在下一
个时刻和“最后”一帧仍然存在的差异，才属于构建。对第一
类请回答 `same`。

规则
- 只比较这些画面所支持的内容。
- 不要假定更新的构建更好，也不要假定一个你无法解释的
  差异就是缺陷。两种倾向都是错误。
- 如果某个行为在这里没有被触发，它的准则就是
  `unobserved`。没看到某样东西，不等于看到它缺失。
- 任何变差的地方都要单独报告，无论其余部分多么微小。
- 判断相对变化。不要给任何一个构建打分。
- 不要提出代码编辑、修复方案或下一步。那不是你的工作。
\end{PV}
\end{promptbox}
